% SPAR Midterm Report Template
% Compatible with arXiv (TeX Live 2023+, pdflatex)
%
% To compile: latexmk -pdf report.tex
% Or use the build scripts: build report.tex

% ============================================================================
% DOCUMENT CLASS
% ============================================================================
% Options: midterm | final | arxiv
% - midterm: For mid-program report
% - final: For end-of-program report
% - arxiv: For converting your report into a preprint for arXiv
\documentclass[midterm]{sparreport}

% ============================================================================
% METADATA
% ============================================================================
\title{Social Context and Judgment Revision in Language Models}

\author{
  \textbf{Hannah Sun}\\
  \texttt{yuerans@uw.edu}\\
  % Optional: Add institution/affiliation (recommended for arXiv mode):
  University of Washington
  % Add more authors (including mentors) with \and:
  \and
  \textbf{Emilio Barkett}\\
  \texttt{eab2291@columbia.edu}\\
  Columbia University
  % Mentor Institution
}

\sparround{Fall 2026}

% ============================================================================
% DOCUMENT BEGINS
% ============================================================================
\begin{document}

\maketitle

% ============================================================================
% ABSTRACT (100-200 words)
% ============================================================================
% Provide an overview of your research question and why it matters for
% addressing risks from advanced AI. Describe your methodological approach
% and any preliminary findings. If working on technical safety or security,
% mention the specific threat model, vector, or vulnerability. Keep this
% accessible to readers outside your subfield. No footnotes or citations.
\begin{abstract}
When AI agents review each other's code or compare research findings, they
must evaluate conflicting judgments. Reliable collaboration requires correcting
errors without abandoning correct answers in response to misleading arguments.
We study how social context affects judgment revision in large language models
(LLMs), varying whether prior answers are described as publicly shared and
whether peer counterarguments use neutral or controlling language.

Using counterfactual replay, we hold questions and counterarguments fixed while
varying contextual framing. Pilots with Llama 3.1 8B, Gemma 3 12B, and Qwen3.5 9B
show that controlling language increases peer-answer adoption in Llama and Qwen,
while public disclosure has little overall effect. The models differ in their
responses to source attribution and their susceptibility to harmful revisions;
Gemma's limited sample does not support a detailed comparison. In re-ask controls,
telling Llama its prior answer was shared increases spontaneous answer changes
relative to describing it as recorded.

Follow-up experiments will test whether attribution and disclosure effects
persist across source identities and alternative prompt formulations.

\end{abstract}

% Keywords (optional but recommended, especially for arXiv submission)
\keywords{AI alignment \and multi-agent systems \and judgment revision \and social influence \and language models}

% ============================================================================
% INTRODUCTION (100-350 words)
% ============================================================================
% Describe the specific problem or question you're investigating and explain
% why it matters. Make clear what gap in existing work you're trying to fill
% or what risk you're addressing. Situate your project relative to existing
% approaches or prior work. Articulate what success would look like for your
% project—what would you need to demonstrate or discover for this research
% to be valuable?
\section{Introduction}
In a team of AI coding agents, one agent might implement a function while another
reviews it and proposes a correction. In a research workflow, agents might
consult different sources and reach conflicting conclusions. In both cases,
an agent must assess whether another agent's reasoning warrants revising its
own judgment. Correcting an error can improve the team's output, but deferring
to a forceful yet misleading argument can introduce new errors.

We examine whether social cues affect this revision process. Our experiments
vary two factors: whether a model's initial answer is described as private or
publicly shared, and whether a conflicting argument uses neutral or controlling
language. The design draws on research on social influence and psychological
reactance \cite{bito2026normative,reynoldstylus2026reactance}, without assuming
that similar behavior in LLMs reflects the same psychological mechanisms.
The question also relates to sycophancy: agreement with another party can come
at the expense of accuracy. Here, we study responses to peer agents rather
than agreement with a user's stated beliefs.

Prior work on LLM conformity and multi-agent debate distinguishes spontaneous
answer changes from peer influence and reasoning-based persuasion
\cite{hao2026flips,hu2026speakerfree}. Source identity and anonymization can also
affect judgment revision \cite{choi2026identity}, and persuasive strategies can
influence model behavior \cite{meincke2026persuading,bozdag2026persuade}.
It remains unclear how directive language interacts with source attribution
and the public disclosure of an earlier answer.

We use controlled counterfactual replay on multi-hop question-answering tasks
to compare responses to the same counterargument under different social cues.
We distinguish beneficial corrections from harmful revisions to assess when
these cues improve accuracy and when they lead models to abandon correct answers.

% ============================================================================
% METHODOLOGY (200-500 words)
% ============================================================================
% Describe your research approach clearly enough that someone could understand
% what you're doing and roughly replicate it. For empirical work, include your
% experimental setup, systems or datasets, and evaluation metrics. For
% theoretical work, outline your framework, key assumptions, and approach to
% arguments or proofs. For policy research, describe your research design,
% data sources, and analytical approach.
\section{Methodology}

% Example subsections (optional):
% \subsection{Experimental Setup}
% \subsection{Dataset and Evaluation}
% \subsection{Analysis Framework}

% Example figure
% \begin{figure}[htbp]
%   \centering
%   \includegraphics[width=0.7\textwidth]{figures/example-figure.pdf}
%   \caption{Example figure caption.}
%   \label{fig:example}
% \end{figure}

% Example table
% \begin{table}[htbp]
%   \centering
%   \caption{Example table caption.}
%   \label{tab:example}
%   \begin{tabular}{lcc}
%     \toprule
%     Method & Accuracy & Time \\
%     \midrule
%     Approach A & 85\% & 2.3s \\
%     Approach B & 92\% & 4.1s \\
%     \bottomrule
%   \end{tabular}
% \end{table}

\subsection{Experimental Framework}
We use counterfactual replay to compare how LLMs respond to conflicting answers
under controlled conversational conditions. Independent agents first answer
the same question without access to one another's responses. We identify
\emph{disagreement events}, where a recipient model and a peer produce different
answers, and replay each event across conditions. The question and substantive
counterargument remain fixed; only the contextual framing changes.

The primary design crosses prior-answer disclosure (private versus public)
with counterargument framing (neutral versus controlling). In public conditions,
the recipient is told that its previous answer was shared with other agents.
Controlling counterarguments direct the recipient to accept the alternative
answer, while neutral counterarguments present the disagreement without that
pressure.

Following work separating spontaneous answer changes from peer-induced revision
\cite{hao2026flips,hu2026speakerfree}, we include controls that remove speaker
attribution, provide only the peer's proposed answer, or ask the recipient to
reconsider without a counterargument. A recorded-answer condition separates
acknowledgment of a prior response from its stated public disclosure. We also
test an independence-oriented persona to examine whether emphasizing independent
judgment changes responses to controlling language.

\subsection{Datasets and Models}
Our initial pilot used ARC-Challenge \cite{clark2018arc} with Llama 3.1 8B
Instruct. We then moved to HotpotQA \cite{yang2018hotpotqa}, where questions
require integrating information across multiple supporting passages. These
tasks allow us to examine disagreements involving multi-step reasoning.

The latest pilots evaluate Llama 3.1 8B, Gemma 3 12B, and Qwen3.5 9B on a common
pool of 1,200 questions balanced across six question categories. Each model
generates its own disagreement events, which we replay under 18 conditions.
After screening and exclusions, the main comparisons include 131 events for
Llama, 18 for Gemma, and 84 for Qwen. Before final exclusions, the runs produced
24,300, 3,240, and 15,120 replay responses, respectively.

\subsection{Evaluation}
Our primary outcome is \emph{answer adoption}: whether the recipient replaces
its initial answer with the peer's proposed answer. We estimate framing and
disclosure effects, including their interaction, through event-level comparisons
with cluster bootstrap confidence intervals.

We separately measure \emph{beneficial corrections}, where an incorrect answer
is replaced by a correct one, and \emph{harmful revisions}, where a correct
answer is replaced by an incorrect alternative. Re-ask controls measure answer
instability in the absence of an external counterargument.

% ============================================================================
% PRELIMINARY RESULTS (200-500 words)
% ============================================================================
% Present what you've found so far. For empirical work, focus on quantitative
% findings with visualizations. For theoretical work, focus on key insights or
% partial results. For policy research, present initial findings from your
% analysis. Use subheadings if helpful. It's normal for midterm results to be
% incomplete.
\section{Preliminary Results}
We initially expected controlling language to produce resistance to a peer's
answer. Instead, it increased adoption in both Llama and Qwen. Public disclosure
had little overall effect on adoption across the three models, providing no
clear evidence that publicly sharing an answer makes a model more resistant
to disagreement. Gemma's estimated framing effect was small (+1.7 percentage
points), but its 18 usable events are insufficient for a detailed behavioral
comparison. Table~\ref{tab:model-comparison} summarizes Llama and Qwen.

\begin{table}[htbp]
  \centering
  \small
  \setlength{\tabcolsep}{5pt}
  \caption{Pilot behavior in Llama and Qwen. Adoption effects are measured in
  percentage points (pp); harmful revision rates are conditional on the initial
  answer being correct. Ranges summarize the reported rates across disclosure
  conditions. Each model generated its own disagreement events, so these are
  descriptive comparisons rather than estimates on a matched event set.}
  \label{tab:model-comparison}
  \begin{tabular}{@{}p{0.40\linewidth}p{0.27\linewidth}p{0.27\linewidth}@{}}
    \toprule
    Measure & Llama 3.1 8B & Qwen3.5 9B \\
    \midrule
    Usable disagreement events & 131 & 84 \\
    Adoption increase under controlling language & +12.8 pp & +16.7 pp \\
    95\% confidence interval & [9.1, 16.4] pp & [11.7, 22.4] pp \\
    Harmful revisions, neutral framing & 71--72\% & 19--20\% \\
    Harmful revisions, controlling framing & 86--89\% & 39--40\% \\
    Effect of naming a peer on the controlling-language effect & Amplifies by 5.3 pp & Reduces by 5.3 pp \\
    Re-ask answer changes: recorded versus shared & 15.5\% versus 27.6\% & Little difference observed \\
    \bottomrule
  \end{tabular}
\end{table}

\subsection{Adoption and Accuracy}
The adoption increase includes harmful revisions, not just successful
corrections. Under neutral framing, Qwen corrected approximately 85--89\% of
initially incorrect answers while abandoning 19--20\% of initially correct
answers. Controlling language raised its harmful revision rate to 39--40\%.
Qwen was more stable than Llama in retaining initially correct answers under
both neutral and controlling framing. Because each model generated its own
disagreement events, this comparison is descriptive rather than evidence of
an intrinsic difference in stability.
Within each model's event set, directive wording therefore increased
susceptibility to misleading counterarguments.

\subsection{Attribution and Baseline Stability}
Naming a peer changed the controlling-language effect in opposite directions:
it amplified the effect in Llama and reduced it in Qwen, relative to presenting
the argument without an attributed speaker. These estimates suggest that
source attribution may alter how models respond to directive wording.

Disclosure produced a separate pattern in Llama's re-ask controls. Without a
counterargument, describing its previous answer as shared increased answer
changes from 15.5\% to 27.6\%, compared with describing it as recorded. Qwen
showed little difference. This result leaves open whether disclosure affects
baseline answer stability even when it has little effect on peer-answer adoption.

% ============================================================================
% CHALLENGES & OPEN QUESTIONS (100-300 words, optional but encouraged)
% ============================================================================
% Be honest about obstacles you've encountered: technical difficulties, resource
% constraints, or conceptual challenges. Note any areas where you're still
% figuring things out or where the path forward is uncertain.
\section{Challenges \& Open Questions}
The main interpretive challenge is separating responses to social context from
ordinary instruction-following. Controlling language explicitly encourages
adoption, so its effect may reflect compliance with a directive rather than
social pressure. Speaker-free and stance-only controls help assess these
explanations. Likewise, the recorded-versus-shared difference may depend on
particular wording rather than public disclosure itself.

The model comparisons also have limits. Each model generates its own
disagreement events, so the selected questions and counterarguments differ.
We cannot attribute the observed differences solely to model behavior without
a matched set of events. Gemma's sample is especially small, partly because of
difficulties obtaining responses in the required format.

Answer correctness is also an incomplete measure of revision quality. A correct
final answer does not establish that the model evaluated the counterargument
appropriately. Because the pilots include many exploratory comparisons,
subgroup patterns need validation in follow-up experiments.

Finally, a behavioral resemblance to human social influence does not establish
a shared psychological mechanism \cite{gupta2026anthropomorphic}. Our follow-ups
will test whether the observed effects persist when we vary the wording and
source of an argument while preserving its substantive content.

% ============================================================================
% NEXT STEPS (150-400 words)
% ============================================================================
% Highlight what you will focus on and what you will try to accomplish during
% the next half of the program. What will you prioritize? What might you
% deprioritize or cut if needed?
\section{Next Steps}
We will prioritize two follow-up experiments based on the attribution and
stability differences in Table~\ref{tab:model-comparison}. Both will retain
answer adoption, beneficial correction, and harmful revision as outcomes.

\subsection{Controlling Language and Source Attribution}
The opposite attribution patterns in Llama and Qwen motivate a closer test of
how source identity interacts with directive language. Prior work shows that
agent identity can affect revision \cite{choi2026identity}, while substantial
updating can occur without an attributed speaker \cite{hu2026speakerfree}.
We will extend the factorial comparison of neutral versus controlling language
and named-peer versus speaker-free arguments, keeping the proposed answer and
reasoning fixed.

We also plan to compare anonymous sources and domain experts to test whether
apparent expertise changes the framing effect. The central question is whether
directive wording has a consistent effect across sources or whether its
influence depends on the authority attributed to the speaker.

\subsection{Recorded Versus Shared Answers}
Llama's higher re-ask instability under shared-answer wording motivates a
separate disclosure experiment. We will compare a plain re-ask baseline with
private, recorded, and shared conditions, using multiple formulations of each
statement. This will test whether the difference persists across wording
changes or depends on specific terms.

We will then introduce neutral and controlling counterarguments while holding
their reasoning fixed, to assess whether disclosure also affects responses to
external disagreement. Work on social context and persuasive framing informs
this design \cite{bito2026normative,meincke2026persuading}, but the immediate
goal is to establish how robust the pilot effect is.

If these effects replicate, we will test additional model families and tasks
beyond HotpotQA. We also plan to evaluate Inspect AI as an evaluation framework
and compare judging approaches for open-ended responses. We will prioritize
replication and prompt variation before expanding to a broader set of tasks.

% ============================================================================
% WORKS CITED / BIBLIOGRAPHY
% ============================================================================
% The reference page lists all sources used for your project, including those
% cited in the review and methods sections. Ensure the citation style is
% consistent throughout.
%
% For arXiv submission: You MUST include the generated .bbl file!
% Compile with: latexmk -pdf report.tex (or pdflatex -> bibtex -> pdflatex -> pdflatex)
% Then submit both .tex and .bbl files to arXiv.

\bibliographystyle{IEEEtran}
\bibliography{references}

\end{document}
